\documentclass[10pt,a4paper]{amsart}
\usepackage[margin=19mm]{geometry}
\usepackage{amsmath,amssymb,mathtools}
\usepackage[scr=rsfs,cal=euler]{mathalfa}
\usepackage{xfrac}
\usepackage[unicode,hidelinks,bookmarks=false]{hyperref}
\newcommand{\ie}{i.e.\ }
\newcommand{\cf}{cf.\ }
\newcommand{\resp}{respectively\ }
\newcommand{\Subsection}{\subsection}
\providecommand{\nobreakdash}{\nobreak\hbox{-}}
\setlength{\emergencystretch}{2em}
\multlinegap0pt
\usepackage{amssymb}
\usepackage[shortlabels]{enumitem}
\newtheorem{proposition}{Proposition}
\title{Geometry of ample/lopsided sets}
\author{Hans--J\"urgen Bandelt, Victor Chepoi, Andreas Dress, Jack Koolen}
\date{Source: arXiv:2603.27835v1 (CC BY 4.0)}
\begin{document}
\maketitle

\noindent\emph{Source and licence.} Geometry of ample/lopsided sets. Version arXiv:2603.27835v1 (CC BY 4.0), distributed by arXiv under the Creative Commons Attribution 4.0 International licence, http://creativecommons.org/licenses/by/4.0/, as recorded in the arXiv licence metadata for this exact version and shown on the abstract page https://arxiv.org/abs/2603.27835v1. Full licence notice: \texttt{licenses/arxiv-2603.27835-CC-BY-4.0.txt}.\par\smallskip

\noindent\textit{Editorial scope.} Counted unit: the article's proposition fiber-connected (source0.tex lines 939--942). The article's complete original proof of it is delivered with it (source0.tex lines 944--986, one \texttt{proof} environment of 43 lines). No mathematics is written, repaired or re-derived: the statement and the proof are the article's own lines, in the article's own notation, and no retained line is substituted or re-flowed.  Retained uncounted as the source-local context the proof needs: lines 912--921.  Nothing was omitted from any retained span, so the package carries no omission record.  No retained line contains a citation, so the wrapper carries no bibliography and no bibliography entry.  No retained line contains a figure environment, an image-inclusion command or drawing code, so no image is delivered and the package declares no asset dependency.  Editorial formatting only, in addition: an AMS \texttt{amsart} preamble that restates the article's own declarations and macros used by the retained lines, namely the environments \texttt{proposition} on separate counters, together with the standard packages those lines need.  The retained environments print in this document as Proposition 1 in order of appearance, whereas the article prints the same items with the numbering of its own full document.  The complete native source of the article as distributed in the arXiv e-print for this exact version is bundled unchanged as \texttt{sources/197386/source0.tex}.
\begin{proposition} \label{fiber-isometric} Let $\Lambda$ be a boundedly compact geodesic space. For a closed subset $K$ of some finite power ${\Lambda}^E$
of $\Lambda$ (endowed with the product metric $d=d^E$), the following statements are equivalent:

\begin{itemize}
\item[\rm(i)] the intrinsic finite-path metric $\delta_{{K},d}$ of $K$ coincides with the restriction $d|_{K}$
of the product distance $d=d^E$ of ${\Lambda}^E$;
\item[\rm(ii)] ${K}$ is a geodesic subspace of ${\Lambda}^E$;
\item[\rm(iii)] ${K}$ is fiber-geodesic.
\end{itemize}
\end{proposition}

\begin{proposition} \label{fiber-connected} Let  $\Lambda$ be a boundedly compact real tree (equipped with the natural metric) in which every segment includes only
finitely many branching points. Then a closed subset $K$ of some finite power ${\Lambda}^E$ of $\Lambda$ is a geodesic subspace of $\Lambda^E$ if and only
if $K$ if fiber-connected.
\end{proposition}

\begin{proof} By Proposition \ref{fiber-isometric}, any closed geodesic subspace of $\Lambda^E$ is fiber-geodesic and hence fiber-connected. Conversely
assume that $K$ is a fiber-connected closed subspace of $\Lambda^E$.  Since $({\Lambda}^E,d)$ is  complete and $K$ is closed, $(K,d|_{K})$ is also
a complete metric space. Thus to prove that $K$ is a geodesic subspace of ${\Lambda}^E$, it suffices
to establish that $K$ is Menger-convex with respect to $d$. We proceed by induction on $\#E.$ Since any proper fiber ${K}_0$ of $K$
is a closed fiber-connected subset of ${\Lambda}^A\times r_0$ for a proper subset $A$ of $E$ and a map $r_0\in \Lambda^{E-A},$
by induction assumption we can assume that ${K}_0$ is a geodesic subspace of ${\Lambda}^A\times \{ r_0\}$ and $\Lambda^E$. Suppose by way
of contradiction that $K$ is not Menger-convex. In view of the induction hypothesis, we may then suppose that
there exist $r_1, r_2 \in {K}$ with
$$[r_1, r_2]_{K}=K \cap [r_1,r_2]_{\Lambda^E}= \{ r_1, r_2\}$$
such that $r_1(x) \neq r_2(x)$ holds for all $x \in E.$

For any $x\in E$, let $\Lambda^x$ be the $x$th factor (a copy of $\Lambda$) of the product ${\Lambda}^E$. Then $\Lambda^{E- \{ x\}}\times r_1(x)$ is the
$(E- \{ x\})$-fiber of $\Lambda^E$ that contains the point $r_1$ of $K$. By $\Lambda^x_+=\Lambda^x_+(r_1,r_2)$  denote the set of all
points $\lambda$ of the real tree
$\Lambda^x$ for which $r_1(x)$ is not between $\lambda$ and $r_2(x),$ that is,
$$\Lambda^x_+:=\{ \lambda\in \Lambda^x: ~ [\lambda,r_1(x)]_{\Lambda^E}\cap [r_1(x),r_2(x)]_{\Lambda^E}\ne \{ r_1(x)\}\}.$$
Then $\Lambda^x_+$ is an open subset of the real tree $\Lambda^x$ and its closure $\overline{\Lambda^x_+}=\Lambda^x_+\cup\{ r_1(x)\}$ is a boundedly
compact subtree of $\Lambda^x$. Trivially, the closure of $\Lambda^{E}_+:=\prod_{x\in E}\Lambda^x_+$ equals
$$\overline{\Lambda^E_+}=\prod_{x\in E} (\Lambda_+^x\cup \{ r_1(x)\}).$$
(Note that $\Lambda^E_+$ resp. $\overline{\Lambda^E_+}$ equal the open resp. closed first orthant of ${\mathbb R}^E$ in the case that $\Lambda={\mathbb R},$
$r_1={\mathbf 0}$ and $r_2>{\mathbf 0}$.)  We claim that
$${K}\cap \overline{\Lambda^{E}_+}=({K}\cap \Lambda^{E}_+)\cup \{ r_1(x)\}.$$
Suppose the contrary, that is, suppose that there exists a point $$r\in ({K}- \{ r_1\})\cap \prod_{y\in E}(\Lambda^y_+\cup \{ r_1(y)\} \mbox{ with } r(x)=r_1(x)$$
for some $x\in E$. Then both $r$ and $r_1$ belong to the fiber $\Lambda^{E- \{ x\}}\times r_1(x)$ and hence are connected by a geodesic $\gamma$ in $K$, according to the induction hypothesis. On the other hand, as $\Lambda$ is a real tree, there exists a point $s$ of $\Lambda^E$ (the median point of $r_1,r,s$) such that
$$[r_1,s]_{\Lambda^E}=[r_1,r_2]_{\Lambda^E}\cap [r_1,r]_{\Lambda^E}.$$
The complement $A=\{ y\in E: ~ r_1(y)\ne r(y)\}$ of the equalizer of $r_1$ and $r$ does not contain the element $x$. Then the geodesic $\gamma$ between $r_1$
and $r$ is included in the $A$-fiber at $r_1$ (and $r$). By the initial hypothesis, $r(y)\in \Lambda^y_+$ and hence $r_1(y)\ne s(y)$ for each $y\in A$. Let
$\epsilon$ be the minimum of the distances of $r_1(y)$ and $s(y)$ for $y\in A$ in the copies of the real tree $\Lambda$. Then the intersection of $\gamma$
with the closed ball of radius $\epsilon$ centered at $r_1$ is included in the box
$$[r_1|_{A},s|_{A}]_{\Lambda^{A}}\times r_1|_{E- A}\subseteq [r_1,r_2]_{\Lambda^E},$$
whence
$$\{ r_1\}\varsubsetneqq\gamma\cap [r_1,r_2]_{\Lambda^E}\subseteq [r_1,r_2]_{K},$$
contrary to the assumption that $[r_1,r_2]_{K}=\{ r_1,r_2\}.$

The intersection ${K}_+:={K}\cap {\Lambda}^E_+$ is an open set in the (topological) subspace $K$ of $\Lambda^E$ that includes $r_2$ but not $r_1$.
We wish to show that ${K}_+$ is closed as well. By what has just been shown,
$$\overline{{K}_+}\subseteq {K}_+\cup \{ r_1\}\subseteq {K}.$$
So suppose that $r_1\in \overline{K_+}.$ Then there exists a sequence $(s_n)$ in $K_+$ converging to $r_1.$ The sequence $(t_n)$ of median points $t_n$ of
$r_1,r_2,$ and $s_n$ $(n\in {\mathbb N})$ also converges to $r_1.$ Since $(s_n)$ is fully included in ${\Lambda}^E_+,$ so is $(t_n).$ It follows from $[r_1,r_2]_{K}=\{ r_1,r_2\}$
that either $t_n(x)=r_2(x)$ or $t_n(x)$ is a branching point of $\Lambda^x$ for each coordinate $x$ and index $n$. This, however, conflicts with the initial requirement
on $\Lambda$ that all segments of $\Lambda$ contain only finitely many branching points. We conclude that $\overline{K_+}={K}_+$ as asserted. Therefore $K_+$ is
both open and closed in $K$, which finally contradicts connectedness of $K$, and the proof is complete.
\end{proof}

\end{document}
